% !TEX program = xelatex
% Biên dịch bằng XeLaTeX hoặc LuaLaTeX để hỗ trợ tiếng Việt qua fontspec/polyglossia
% (compile với: xelatex main.tex — chạy 2 lần để mục lục/tham chiếu chương đúng)
\documentclass[a4paper,12pt]{report}

\usepackage{fontspec}
\usepackage{polyglossia}
\setmainlanguage{vietnamese}
\setmainfont{Times New Roman} % đổi sang font có sẵn trên máy nếu không có (vd: Noto Serif)

\usepackage[margin=2.5cm]{geometry}
\usepackage{setspace}
\onehalfspacing

\usepackage{graphicx}
\usepackage{booktabs}
\usepackage{longtable}
\usepackage{hyperref}
\usepackage{amsmath}
\usepackage{listings}
\usepackage{xcolor}

\lstset{
  basicstyle=\ttfamily\small,
  breaklines=true,
  frame=single,
  columns=fullflexible,
}

\hypersetup{
  colorlinks=true,
  linkcolor=blue,
  citecolor=blue,
  urlcolor=blue,
}

\title{Hệ thống chấm điểm bài tập lập trình tự động\\ dựa trên mô hình ngôn ngữ lớn dưới 7 tỷ tham số}
\author{Khang Nguyen}
\date{\today}

\begin{document}

\maketitle
\tableofcontents

% ============================================================
\chapter{Tổng quan tài liệu}
\label{ch:lit-review}

\section{Phạm vi}

Đề tài này nằm ở giao điểm của hai hướng nghiên cứu: chấm điểm tự động bài
tập lập trình bằng LLM, và hướng nghiên cứu rộng hơn về LLM-as-a-judge —
làm sao để lấy được một điểm số đáng tin cậy từ một mô hình ngôn ngữ. Các
mục dưới đây được tổ chức theo từng quyết định thiết kế trong
\texttt{docs/architecture.md}: mỗi mục nêu quyết định thiết kế, tài liệu
tham khảo làm căn cứ, và — ở những chỗ có thực nghiệm trong repo kiểm chứng
thực nghiệm cho luận điểm đó — dẫn chiếu tới kết quả tương ứng
(\texttt{docs/thesis/results.md}).

\section{Từ chấm điểm toàn thể (holistic) đến đánh giá có cấu trúc, phân rã}

Vấn đề gốc của cách làm ``cứ hỏi thẳng LLM cho điểm'' là một lệnh gọi tự do
(free-form) tạo ra một điểm số toàn thể duy nhất có phương sai (variance)
rất cao giữa các lần chạy — mô hình suy luận khác nhau ở mỗi lần gọi, kể cả
khi giải mã tham lam (greedy decoding), vì văn bản tự do cho mô hình nhiều
đường suy luận khác nhau nhưng đều hợp lý để đi đến cùng một con số cuối
cùng. Bài khảo sát về LLM-as-a-Judge (arXiv:2411.15594
\cite{survey_llm_judge}) và công trình về đánh giá theo rubric có cấu trúc
(\textit{From Holistic Evaluation to Structured Criteria: Rubrics Across
the Evolving LLM Landscape}, arXiv:2606.08625
\cite{rubric_decomposition}) đều hội tụ về cùng một giải pháp: phân rã một
phán đoán toàn thể thành các kiểm tra (check) nhỏ, độc lập, có cấu trúc, và
tổng hợp kết quả bên ngoài mô hình thay vì để mô hình tự tạo ra hoặc tự
điều chỉnh con số cuối cùng. Đây chính là động lực trực tiếp cho hình dạng
pipeline của đề tài — chấm điểm đúng/sai (correctness) bằng thực thi trong
sandbox cộng với một lệnh gọi LLM hẹp, có cấu trúc theo schema, cho mỗi
tiêu chí rubric còn lại (\texttt{src/grader/rubric\_checks.py}), tổng hợp
bằng một tổng có trọng số cố định trong code thuần
(\texttt{src/grader/aggregator.py}) — được mô tả đầy đủ trong
\texttt{docs/architecture.md}.

Một vấn đề liên quan nhưng khác biệt được nêu trong \textit{Reliability
without Validity: Systematic Evaluation of LLM-as-a-Judge}
(arXiv:2606.19544 \cite{reliability_without_validity}): một ``giám khảo''
(judge) có thể đáng tin cậy (reliable — nhất quán) mà không nhất thiết hợp
lệ (valid — đúng, hoặc khớp với một chuẩn con người thực sự). Sự phân biệt
này ánh xạ trực tiếp vào hai nhánh thực nghiệm riêng biệt của đề tài — Pha
1 đo độ tin cậy (tỷ lệ khớp tuyệt đối — exact-match — qua nhiều lần chạy
lặp lại, mục ``Consistency'' trong \texttt{docs/thesis/results.md}) và Pha
2 đo một đại lượng đại diện cho tính hợp lệ (độ tương quan Spearman với
nhãn tham chiếu của bộ dữ liệu, mục ``Ablation'' trong
\texttt{docs/thesis/results.md}) — và cũng là lý do đề tài luôn gọi cẩn
thận các con số tương quan ở Pha 2 là ``mức độ đồng thuận với nhãn dựa trên
luật'' (agreement with rule-based labels), chứ không phải ``mức độ đồng
thuận với giám khảo con người'': nhãn tham chiếu cho \texttt{efficiency}
và \texttt{code\_style} do một người duy nhất (chính tác giả đề tài) tự gán
nhãn, không phải nhãn chuẩn độc lập từ con người
(\texttt{data/script.md}). Đáng tin cậy mà không hợp lệ là một cái bẫy dễ
gặp ở đây — một pipeline luôn trả về cùng một điểm số \emph{sai} ở mọi lần
chạy vẫn được coi là ``đáng tin cậy'' nhưng vô dụng — vì vậy cả hai nhánh
thực nghiệm đều được coi là cần thiết, không thể thay thế cho nhau.

\section{Chấm điểm từng bước mà không có phương sai của văn bản tự do}

Hai bài báo chuyên về chấm điểm bài tập lập trình, CodEv (arXiv:2501.10421
\cite{codev}) và StepGrade (\textit{Grading Programming Assignments with
Context-Aware LLMs}, arXiv:2503.20851 \cite{stepgrade}), đều báo cáo rằng
chấm điểm theo kiểu chuỗi suy luận (chain-of-thought) hoặc từng bước
(step-by-step) cho kết quả tốt hơn một lệnh gọi toàn thể duy nhất. Đây là
một mâu thuẫn thực sự với yêu cầu về tính nhất quán ở trên: suy luận
chain-of-thought bằng văn bản tự do chính là loại sinh văn bản không ràng
buộc mà thực nghiệm về tính nhất quán ở Pha 1 được thiết kế để nhạy cảm với
nó. Giải pháp của đề tài này là lấy lợi ích của ``từng bước'' mà hai bài
báo trên báo cáo, nhưng không phải trả giá bằng phương sai của văn bản tự
do — việc phân rã diễn ra dưới dạng các \emph{lệnh gọi} mô hình riêng biệt,
mỗi lệnh gọi giới hạn trong đúng một tiêu chí rubric với đầu ra dạng đóng
(enum hoặc boolean) nhỏ gọn thông qua giải mã ràng buộc theo văn phạm
(grammar-constrained decoding), thay vì một chuỗi suy luận dài hướng tới
một con số. Phần mô tả pipeline trong \texttt{docs/architecture.md} nêu rõ
đây là câu trả lời của đề tài cho phát hiện của CodEv/StepGrade.

\section{Đặt correctness bên ngoài mô hình}

Hai bài báo khác củng cố việc giữ chấm điểm đúng/sai (correctness) hoàn
toàn bên ngoài LLM. \textit{Automating Autograding: LLMs as Test Suite
Generators} (arXiv:2411.09261 \cite{autograding_test_suites}) và
\textit{Multi-Programming Language Sandbox for LLMs} (arXiv:2410.23074
\cite{multi_lang_sandbox}) đều làm việc với công cụ liên quan đến LLM xoay
quanh việc thực thi code thật, thay vì để LLM đánh giá correctness bằng
cách đọc code. Đề tài này giữ lập trường tương tự, vì lý do tính xác định
(determinism) cũng nhiều như vì lý do độ chính xác: correctness được chấm
hoàn toàn bằng cách chạy code nộp bài với các test case thật trong sandbox
(\texttt{src/grader/sandbox.py}) — có tính xác định theo cấu trúc, không
có LLM tham gia, do đó không có bề mặt phương sai nào trên tiêu chí rubric
này. Chính ghi chú xây dựng bộ dữ liệu (\texttt{data/script.md}, mục
``Notable finding from running the real pipeline'') cũng độc lập củng cố
một lưu ý ngầm trong hướng nghiên cứu này: correctness từ sandbox chỉ tốt
bằng bộ test case đứng sau nó — hai bài nộp được thiết kế cố tình có lỗi
vẫn đạt 100\% test vì bộ test không bao phủ đúng lỗi đó, một lời nhắc rằng
``ground truth từ thực thi'' không tự động là ground truth đầy đủ.

\section{Tính xác định dưới batched serving: một khoảng trống trong tài liệu}

Câu hỏi thực nghiệm khó hơn của yêu cầu về tính nhất quán — liệu dynamic
batching trong một serving engine có thể phá vỡ tính xác định bit-chính-xác
ngay cả ở temperature = 0, do tính không kết hợp (non-associativity) của
số học dấu phẩy động qua các tổ hợp batch khác nhau — được
\texttt{docs/architecture.md} coi là một hiện tượng đã được ghi nhận, đo
đạc trong lĩnh vực serving LLM nói chung, là động lực cho thực nghiệm về
tính nhất quán ở Pha 1. Tìm kiếm trong các tài liệu đã trích dẫn ở trên
không thấy bài báo nào đo trực tiếp tỷ lệ khớp tuyệt đối (exact-match) của
việc \emph{chấm điểm} bằng LLM dưới điều kiện batched serving;
\texttt{docs/thesis/results.md} ghi nhận rõ đây là một khoảng trống mà
chính thực nghiệm về tính nhất quán của đề tài (Pha 1, chạy trên
continuous-batching của llama.cpp) phần nào giải quyết được, kèm lưu ý về
phạm vi rằng một serving engine và một họ mô hình duy nhất ở một mức độ
concurrency không thiết lập được kết luận tổng quát. Khoảng trống này là
một cách đóng khung khả dĩ cho đóng góp của đề tài: một phép đo tính nhất
quán thực nghiệm cho một pipeline chấm điểm hybrid-decomposed cụ thể, dưới
một serving engine cụ thể, thay vì một tuyên bố tổng quát về tính xác định
của serving LLM.

\section{Tóm tắt: thiết kế của đề tài bám theo tài liệu như thế nào}

\begin{table}[h]
\centering
\begin{tabular}{p{4cm}p{4cm}p{5cm}}
\toprule
\textbf{Quyết định thiết kế} & \textbf{Tài liệu căn cứ} & \textbf{Nơi triển khai} \\
\midrule
Phân rã rubric thành các check độc lập, có cấu trúc &
  \cite{survey_llm_judge}, \cite{rubric_decomposition} &
  \texttt{src/grader/rubric\_checks.py}, \texttt{src/grader/schemas.py} \\
Tách riêng độ tin cậy và tính hợp lệ như hai luận điểm không thay thế
  được cho nhau &
  \cite{reliability_without_validity} &
  Pha 1 (consistency) và Pha 2 (ablation agreement),
  \texttt{docs/thesis/results.md} \\
Lấy lợi ích chấm điểm từng bước mà không có phương sai của CoT tự do &
  \cite{codev}, \cite{stepgrade} &
  Các lệnh gọi riêng cho từng tiêu chí thay vì một chuỗi CoT \\
Chấm correctness bằng thực thi thật, không phải phán đoán của LLM &
  \cite{autograding_test_suites}, \cite{multi_lang_sandbox} &
  \texttt{src/grader/sandbox.py} \\
Đo tính xác định của batched serving bằng thực nghiệm thay vì giả định &
  Khoảng trống đã xác định, chưa tìm thấy công trình trực tiếp trước đó &
  \texttt{experiments/consistency\_test.py}, Pha 1 \\
\bottomrule
\end{tabular}
\caption{Ánh xạ giữa quyết định thiết kế và tài liệu tham khảo}
\end{table}

\section{Lưu ý về mức độ đầy đủ (dành cho người đọc/người hướng dẫn)}

Bản dự thảo này chỉ nêu những luận điểm đã có căn cứ sẵn trong mô tả một
dòng của \texttt{docs/architecture.md} cho mỗi bài báo được trích dẫn.
Chương này \emph{không} thêm chi tiết phương pháp, số liệu báo cáo, hay
phát hiện cụ thể của các bài báo đó ngoài những gì đã nêu ở đó, vì repo
hiện chưa có toàn văn hay tóm tắt (abstract) của các bài báo để đối chiếu
xa hơn — muốn thêm chi tiết hơn (ví dụ: số liệu accuracy CodEv/StepGrade tự
báo cáo, hoặc cơ chế batching cụ thể đứng sau luận điểm về tính không kết
hợp của số học dấu phẩy động) cần bổ sung nguồn đó vào danh sách tham khảo
của \texttt{docs/architecture.md} trước, để đường dẫn trích dẫn vẫn có thể
kiểm chứng được. Ngoài ra, danh mục tài liệu tham khảo ở cuối file này hiện
\textbf{chưa có tên tác giả/năm xuất bản/hội nghị đầy đủ} cho 7 trích dẫn
trên — cần tra lại từng arXiv ID trên arxiv.org và điền đầy đủ trước khi
nộp luận văn chính thức.

% ============================================================
\chapter{Phương pháp đề xuất}
\label{ch:proposed-method}

\section{Tổng quan kiến trúc}

Hệ thống chấm điểm là một pipeline tuần tự duy nhất
(\texttt{src/grader/pipeline.py}), \emph{không} phải một khung điều phối
đa tác tử (multi-agent orchestration framework). Với một mẫu bài nộp gồm đề
bài, rubric, code nộp bài và tập test case, pipeline thực hiện bốn bước:

\begin{enumerate}
  \item \textbf{Chấm correctness bằng thực thi thật} trong sandbox
    (\texttt{sandbox.run\_submission()}) — có tính xác định theo cấu
    trúc, không có LLM tham gia.
  \item \textbf{Ba lệnh gọi LLM độc lập, có cấu trúc}, mỗi lệnh gọi chấm
    đúng một tiêu chí rubric còn lại: \texttt{efficiency},
    \texttt{code\_style}, \texttt{edge\_case\_handling}
    (\texttt{rubric\_checks.py}), ở nhiệt độ (temperature) = 0, seed cố
    định.
  \item \textbf{Tổng hợp điểm cuối cùng bằng code thuần}
    (\texttt{aggregator.aggregate()}) — một tổng có trọng số cố định,
    không bao giờ do LLM tính hoặc điều chỉnh.
  \item (Tùy chọn, không ảnh hưởng điểm số) \textbf{Sinh phản hồi giải
    thích} theo dòng code cụ thể (\texttt{feedback.py}), chạy sau khi đã
    có kết quả chấm, tách biệt hoàn toàn khỏi quá trình tính điểm.
\end{enumerate}

Sơ đồ luồng dữ liệu (theo \texttt{docs/architecture.md}):

\begin{verbatim}
input: statement, rubric, submission_code, test_cases
   |
   +-- sandbox.run_submission()          -> correctness (xác định, không LLM)
   |
   +-- rubric_checks.check_efficiency()          -> JSON có cấu trúc, temp=0, seed cố định
   +-- rubric_checks.check_code_style()          -> JSON có cấu trúc, temp=0, seed cố định
   +-- rubric_checks.check_edge_case_handling()  -> JSON có cấu trúc, temp=0, seed cố định
   |
   +-- aggregator.aggregate()             -> final_score = sum(sub_score*weight)/sum(weight)*100
\end{verbatim}

\section{Vì sao không dùng agent framework}

Quyết định không dùng một khung điều phối đa tác tử (LangChain, CrewAI,
v.v.) xuất phát trực tiếp từ yêu cầu cứng của đề tài — tính nhất quán
(consistency) giữa các lần chạy lặp lại trên cùng một đầu vào. Mỗi lớp điều
phối thêm vào (lập kế hoạch từng bước, định tuyến gọi công cụ, giao tiếp
giữa các tác tử) là một nơi có thể làm rò rỉ thêm phi-xác-định
(non-determinism) hoặc token, mà không mang lại lợi ích gì ở đây: không có
khám phá nhiều bước hay sử dụng công cụ (tool use) cần điều phối — chỉ có
ba lệnh gọi phân loại độc lập và một phép cộng ở phía code. Việc phân rã
rubric thành các kiểm tra nhỏ, độc lập, có cấu trúc — như tài liệu
LLM-as-a-judge đã chỉ ra (Chương~\ref{ch:lit-review}) — là một bài toán
thiết kế dữ liệu/prompt, không phải bài toán điều phối.

\section{Chấm correctness: thực thi trong sandbox}

Correctness được chấm bởi \texttt{sandbox.run\_submission()}: nạp code nộp
bài vào một namespace Python riêng bằng \texttt{exec()}, sau đó chạy hàm
tương ứng (\texttt{func\_name}) lần lượt trên từng test case
(\texttt{(args, expected)}), so sánh kết quả thực tế với kết quả kỳ vọng.
Hàm trả về một cấu trúc kết quả gồm \texttt{pass\_count}, \texttt{total},
\texttt{pass\_rate}, danh sách lỗi cụ thể (\texttt{errors}) và một
\texttt{status} (\texttt{ran}, \texttt{compile\_error}, hoặc
\texttt{func\_not\_found} nếu code không biên dịch được hoặc không định
nghĩa đúng hàm cần chấm). Điểm correctness dưới dạng số
(\texttt{score\_correctness()}) chính là \texttt{pass\_rate} — tỷ lệ test
case vượt qua, không qua bất kỳ bước diễn giải nào của LLM. Đây là logic
được dùng chung với \texttt{data/build\_dataset.py} khi xây bộ dữ liệu, để
pipeline chấm điểm thực tế và bộ dữ liệu tham chiếu không lệch nhau theo
thời gian.

Một giới hạn quan trọng của cách tiếp cận này, đã được ghi nhận thực nghiệm
trong \texttt{data/script.md}: chất lượng của correctness phụ thuộc trực
tiếp vào độ bao phủ của bộ test case — nếu bộ test yếu, sandbox vẫn chạy
đúng về mặt kỹ thuật nhưng cho ra điểm sai (hai bài nộp cố tình có lỗi vẫn
đạt 100\% vì test không bao phủ đúng lỗi đó; xem
Chương~\ref{ch:lit-review}, mục về đặt correctness ngoài mô hình).

\section{Ba kiểm tra rubric có cấu trúc, độc lập}

Với mỗi tiêu chí rubric ngoài correctness, pipeline gọi đúng một prompt hẹp
(\texttt{prompts/*.txt}), yêu cầu mô hình trả về một đối tượng JSON nhỏ
khớp với một schema Pydantic cố định (\texttt{schemas.py}), thông qua giải
mã có ràng buộc (\texttt{response\_format: json\_schema}, dựa trên bộ máy
văn phạm GBNF của llama.cpp). Nhờ vậy không có con số tự do nào để có thể
thiếu nhất quán — chỉ có một enum hoặc boolean bị giới hạn, sau đó được ánh
xạ sang điểm 0..1 theo một công thức cố định ở phía code.

\subsection{Efficiency}

\texttt{EfficiencyCheck} yêu cầu mô hình trả về \texttt{meets\_expected\_}
\texttt{complexity} (boolean) và \texttt{observed\_pattern} (một cụm từ
ngắn mô tả mẫu độ phức tạp quan sát được, ví dụ \texttt{"nested loop
O(n\textasciicircum 2)"}). Điểm số:
\[
\text{score}_{\text{efficiency}} =
\begin{cases}
1.0 & \text{nếu } \texttt{meets\_expected\_complexity} = \text{true} \\
0.4 & \text{ngược lại}
\end{cases}
\]
Trường \texttt{observed\_pattern} không ảnh hưởng điểm số — nó được giữ lại
để phục vụ sinh phản hồi giải thích (Mục~\ref{sec:feedback}), thay vì bị bỏ
đi hoàn toàn như một phiên bản trước đó của hàm này từng làm.

\subsection{Code style}

\texttt{StyleCheck} yêu cầu mô hình trả về đúng một giá trị enum thứ bậc
\texttt{OrdinalLevel} $\in \{\texttt{poor}, \texttt{medium},
\texttt{clean}\}$, ánh xạ sang điểm theo bảng cố định:

\begin{table}[h]
\centering
\begin{tabular}{ll}
\toprule
\textbf{Mức (level)} & \textbf{Điểm} \\
\midrule
\texttt{clean} & 1.0 \\
\texttt{medium} & 0.6 \\
\texttt{poor} & 0.2 \\
\bottomrule
\end{tabular}
\caption{Ánh xạ mức code style sang điểm 0..1}
\end{table}

\subsection{Edge-case handling}

\texttt{EdgeCaseCheck} yêu cầu mô hình trả về \texttt{handled} (boolean) và
\texttt{missing\_cases} (danh sách chuỗi, có thể rỗng, liệt kê các trường
hợp biên chưa được xử lý). Công thức điểm:
\[
\text{score}_{\text{edge\_case}} =
\begin{cases}
1.0 & \text{nếu } \texttt{handled} = \text{true} \\
\max\bigl(0.25,\; 1.0 - 0.25 \times \max(1, |\texttt{missing\_cases}|)\bigr)
  & \text{nếu } \texttt{handled} = \text{false}
\end{cases}
\]
Việc lấy $\max(1, \cdot)$ trong nhánh \texttt{false} là một điểm sửa lỗi có
chủ đích: nếu mô hình trả về \texttt{handled=False} nhưng
\texttt{missing\_cases} rỗng (không nêu được trường hợp cụ thể), công thức
cũ sẽ tính $1.0 - 0.25 \times 0 = 1.0$ — tức là âm thầm ghi đè phán quyết
``không xử lý được'' của chính mô hình thành điểm tuyệt đối. Danh sách rỗng
ở đây có nghĩa là ``chưa xác định cụ thể'', không phải ``không bị trừ
điểm'', nên điểm sàn được giữ ở 0.25 thay vì cho phép điểm tối đa.

\section{Tổng hợp điểm cuối cùng}

Điểm cuối cùng được tính hoàn toàn ở phía code, không qua LLM
(\texttt{aggregator.aggregate()}):
\[
\text{final\_score} = \text{round}\!\left(
  \frac{\sum_{k} \text{sub\_score}_k \times \text{weight}_k}
       {\sum_{k} \text{weight}_k} \times 100,\; 2 \right)
\]
với $k \in \{\texttt{correctness}, \texttt{efficiency},
\texttt{code\_style}, \texttt{edge\_case\_handling}\}$, và các trọng số
\texttt{weight}$_k$ lấy trực tiếp từ rubric của từng mẫu bài (tổng các
trọng số không bắt buộc phải là 100, công thức tự chuẩn hóa qua phép chia
cho $\sum_k \text{weight}_k$).

\section{Sinh phản hồi giải thích (không ảnh hưởng điểm số)}
\label{sec:feedback}

Sau khi đã có kết quả chấm, \texttt{feedback.py} có thể sinh phản hồi kiểu
giáo viên rà bài, dựa trên \texttt{CodeFeedback} — một tóm tắt 1--2 câu
cộng một danh sách \texttt{FeedbackItem} (mỗi mục gắn với một số dòng cụ
thể trong code, nêu vấn đề, giải thích ngắn, và đề xuất sửa). Bước này chạy
\emph{sau} và \emph{tách biệt} khỏi quá trình tính điểm — nó dùng lại kết
quả sandbox/rubric đã có để giải thích điểm số, nhưng không có khả năng làm
thay đổi điểm số đó.

\section{Lựa chọn mô hình và cấu hình phục vụ (serving)}

Mô hình mặc định của đề tài là \textbf{Qwen2.5-Coder-7B-Instruct} (định
dạng GGUF, lượng tử hóa Q4\_K\_M) — mô hình lớn nhất còn nằm trong phạm vi
dưới 7 tỷ tham số mà đề tài đặt ra. Ban đầu đề tài dự kiến dùng vLLM, nhưng
máy phát triển là Apple Silicon (M1 Pro, không có CUDA) — vLLM không có
backend GPU được hỗ trợ trên nền tảng này, nên đề tài chuyển sang
\texttt{llama-server} của llama.cpp: tương thích API kiểu OpenAI, tăng tốc
bằng Metal, và hỗ trợ giải mã ràng buộc theo văn phạm GBNF cho đầu ra JSON.
Cấu hình xác định (\texttt{configs/model.yaml}) cố định
\texttt{temperature=0}, \texttt{seed=0}, và
\texttt{guided\_decoding\_backend=gbnf}.

Vì \texttt{temperature=0} + seed cố định + giải mã ràng buộc \emph{không}
tự động đảm bảo lặp lại chính xác từng bit dưới một serving engine có
batching (do tính không kết hợp của số học dấu phẩy động qua các tổ hợp
batch khác nhau), đề tài coi đây là điều cần đo thực nghiệm chứ không phải
điều một cờ cấu hình (config flag) đảm bảo sẵn —
\texttt{experiments/consistency\_test.py} là phép kiểm định thực nghiệm
tương ứng (chi tiết và kết quả: \texttt{docs/thesis/results.md}, ngoài
phạm vi chương Phương pháp này).

\section{Thiết kế bộ dữ liệu thí điểm}

Bộ dữ liệu thí điểm (\texttt{data/dataset.json}, sinh từ
\texttt{data/problems.py} qua \texttt{data/build\_dataset.py}) gồm 65 mẫu
trên 16 bài toán gốc, mỗi bài toán có 4--5 lời giải mô phỏng các loại lỗi
khác nhau (đúng-sạch, đúng-nhưng-kém-hiệu-quả/kém-style, và vài dạng sai
khác nhau). Mỗi mẫu gồm: đề bài, rubric, code nộp bài, kết quả thực thi
(\texttt{pass\_count}, \texttt{total\_tests}, \texttt{pass\_rate},
\texttt{status}, \texttt{sample\_errors}), các điểm thành phần
(\texttt{sub\_scores\_0\_to\_1}), điểm cuối cùng, và nguồn gốc của từng
nhãn (\texttt{annotation\_source}).

\begin{table}[h]
\centering
\begin{tabular}{lll}
\toprule
\textbf{Tiêu chí} & \textbf{Nguồn} & \textbf{Độ tin cậy} \\
\midrule
correctness & Chạy code thật với test case trong sandbox & Khách quan hoàn toàn \\
edge\_case\_handling & Suy ra trực tiếp từ pass\_rate thực thi & Khách quan hoàn toàn \\
efficiency & Quy tắc cố định, tự phân loại độ phức tạp thuật toán
  khi viết bài nộp & Một người tự gán nhãn, KHÔNG phải ground truth thật \\
code\_style & Quy tắc 3 mức cố định (clean/medium/poor), tự gán & Một người tự gán nhãn,
  KHÔNG phải ground truth thật \\
\bottomrule
\end{tabular}
\caption{Nguồn gốc và độ tin cậy của từng nhãn trong bộ dữ liệu thí điểm
(\texttt{data/script.md})}
\end{table}

\textbf{Hạn chế quan trọng nhất, cần nhắc lại xuyên suốt luận văn:} dữ liệu
này là code do chính tác giả viết với lỗi cố ý chèn vào, tự phân loại —
không phải bài làm của sinh viên thật, và không có điểm số từ giảng viên
thật. Do đó bộ dữ liệu này \emph{có thể} dùng để kiểm tra pipeline chạy
đúng và đo tính nhất quán của LLM (không cần ground truth thật, chỉ cần đầu
vào cố định), nhưng \emph{không thể} dùng để báo cáo ``mô hình đạt X\% độ
tương quan với con người'' trong luận văn — vì ``con người'' ở đây là chính
tác giả tự gán nhãn theo luật, không phải giảng viên chấm độc lập. Bước
tiếp theo cần thiết (T3.2 trong \texttt{docs/roadmap.md}, hiện đang tạm
hoãn theo quyết định của người dùng) là nhờ 1--2 giảng viên/trợ giảng chấm
độc lập một mẫu con để tính hệ số Cohen's kappa / ICC giữa nhãn theo luật
và nhãn con người thật.

\section{Tóm tắt đóng góp phương pháp luận}

Ba điểm đóng góp chính của phương pháp, tổng hợp lại:

\begin{enumerate}
  \item Phân rã rubric thành các kiểm tra nhỏ, độc lập, có đầu ra dạng
    đóng (schema-constrained) thay vì một phán đoán toàn thể tự do —
    được thực nghiệm hóa qua so sánh 3 nhánh (free-form, hybrid-freeform,
    hybrid-decomposed) ở Pha 2.
  \item Tách hoàn toàn correctness (thực thi thật, xác định) khỏi các tiêu
    chí còn lại (LLM, có cấu trúc) — loại bỏ LLM khỏi con đường chấm điểm
    quan trọng nhất về mặt xác định.
  \item Đo tính nhất quán giữa các lần chạy như một đại lượng thực nghiệm
    độc lập với tính hợp lệ (đồng thuận với nhãn tham chiếu), thay vì giả
    định rằng cấu hình \texttt{temperature=0} + seed cố định là đủ để đảm
    bảo tính xác định dưới serving có batching.
\end{enumerate}

% ============================================================
% Các chương còn lại (Thực nghiệm, Kết quả, Hạn chế) sẽ được thêm vào khi
% các mục T5.3, T5.4 trong docs/roadmap.md hoàn thành.

% ============================================================
% Tài liệu tham khảo — inline thebibliography (không cần bibtex/biber).
% LƯU Ý: chỉ có tiêu đề + arXiv ID, chưa có tên tác giả/năm/hội nghị đầy đủ
% (không có sẵn trong repo) — cần tra lại trên arxiv.org trước khi nộp.
\begin{thebibliography}{9}

\bibitem{survey_llm_judge}
Survey on LLM-as-a-Judge. arXiv:2411.15594.

\bibitem{rubric_decomposition}
From Holistic Evaluation to Structured Criteria: Rubrics Across the
Evolving LLM Landscape. arXiv:2606.08625.

\bibitem{reliability_without_validity}
Reliability without Validity: Systematic Evaluation of LLM-as-a-Judge.
arXiv:2606.19544.

\bibitem{codev}
CodEv: Automated Grading Framework. arXiv:2501.10421.

\bibitem{stepgrade}
StepGrade: Grading Programming Assignments with Context-Aware LLMs.
arXiv:2503.20851.

\bibitem{autograding_test_suites}
Automating Autograding: LLMs as Test Suite Generators. arXiv:2411.09261.

\bibitem{multi_lang_sandbox}
Multi-Programming Language Sandbox for LLMs. arXiv:2410.23074.

\end{thebibliography}

\end{document}
